% Draft Methods --- analytical half, rewritten with inline two-register (plain +
% formal) and the new analyses (Shapley, Moran, mapping-scale, 6-layer biology,
% 500-tree block sweep, ESM structural comparison). Staging file (2026-06-29).
%
% CARRIED OVER from current manuscript.tex with only minor updates (NOT reproduced
% here): Conceptual framework (tau definition, Eq. 1), Soil data compilation
% (OpenLandMap compendium), Data QC, Bulk-density gap-filling, Mean transit time
% calculation, Predictor variables (Climate/Edaphic/Land use). The Biological
% predictor subsection gains the MAIN/sensitivity note below.

\section{Methods}
\label{sec:methods}

% ---------------------------------------------------------------------------
\subsection{Modelling design and the attribution problem}
% ---------------------------------------------------------------------------
We organise predictors into four mechanistically distinct groups --- Climate,
Edaphic, Land use, and Biological (Table~\ref{tab:predictors}) --- and model
$\ln\tau$ with Random Forests. Productivity variables (NPP, leaf area index) are
excluded by design: because $\tau$ is computed as SOC stock divided by an
NPP-derived input (Eq.~\ref{eq:mtt}), admitting productivity as a predictor would
make the model partly circular.

The central analytical question is how much of the predictable variation in $\tau$
belongs to each group. This is not a trivial accounting exercise, because the
groups are spatially correlated --- climate is itself a soil-forming factor, so
climate and edaphic descriptors carry overlapping information (Sec.~\ref{sec:intro})
--- and correlated predictors cannot be credited by simply adding them one at a
time, since the credit each receives then depends on the order of addition. We
therefore attribute variance with an order-independent decomposition (Shapley
values, Sec.~\ref{sec:shapley}) rather than with nested increments, and we test
separately whether the variance left unexplained by climate is spatially organised
(Sec.~\ref{sec:moran}), where it acts (Sec.~\ref{sec:mapscale}), and how each
driver acts (Sec.~\ref{sec:ale}).

% ---------------------------------------------------------------------------
\subsection{Random Forest fitting and spatial cross-validation}
\label{sec:rf}
% ---------------------------------------------------------------------------
All models were fitted with the \texttt{ranger} package \citep{Wright2017}
(500~trees, minimum node size~5, $m_\text{try}=\lfloor\sqrt{p}\rfloor$), on
$\ln\tau$, with predictions back-transformed by exponentiation. Skill was
estimated by \emph{spatial block} cross-validation: the globe was partitioned into
$2\textdegree\times2\textdegree$ blocks ($\sim$200~km), blocks were assigned to ten
folds, and each fold's blocks were predicted from a model trained on the other
nine. Holding out whole blocks --- rather than individual points --- prevents a
site near a fold boundary from being predicted by an almost-coincident training
neighbour, so the reported $R^2$ measures generalisation to unseen regions rather
than interpolation between adjacent samples.

The $2\textdegree$ block exceeds the spatial range of the beyond-climate residual
(empirical semivariogram range $\approx16$~km; Sec.~\ref{sec:mapscale}), so
held-out blocks are effectively independent of the training data. The attribution
is also insensitive to the block size: although the absolute cross-validated $R^2$
declines as blocks enlarge (stricter spatial separation, fewer effective degrees
of freedom), repeating the decomposition with $1\textdegree$ and $5\textdegree$
blocks (at the production 500 trees) shifts each group's share of attributable
variance by at most 2.5 percentage points, leaving the ranking and the
climate-versus-non-climate split unchanged.

% ---------------------------------------------------------------------------
\subsection{Attribution of variance among process classes (Shapley)}
\label{sec:shapley}
% ---------------------------------------------------------------------------
\textit{In brief.} To state how much of the predictable variation in $\tau$ each
process class accounts for, we borrow a tool from cooperative game theory. When
players (here, predictor groups) contribute partly overlapping value, the fair way
to credit each is to average its marginal contribution over every possible order in
which players could join the team. Applied to explained variance, this Shapley
share is independent of any arbitrary ordering --- the property that nested,
add-one-at-a-time increments lack when predictors are correlated.

\textit{Formally.} Let $N$ be the four groups and $v(S)$ the spatially
cross-validated $R^2$ of the Random Forest using exactly the groups in coalition
$S\subseteq N$. We fit all $2^{4}=16$ coalitions on a single common set of rows ---
those complete across all groups ($n=154{,}675$) --- so that every $v(S)$ is
estimated on the same observations and the coalition scores are mutually
comparable. The Shapley value of group $i$ is
\begin{equation}
  \phi_i \;=\!\! \sum_{S\subseteq N\setminus\{i\}}
  \frac{|S|!\,\bigl(|N|-|S|-1\bigr)!}{|N|!}\,
  \bigl[\,v(S\cup\{i\})-v(S)\,\bigr],
  \label{eq:shapley}
\end{equation}
the order-averaged marginal contribution of $i$; the $\phi_i$ sum exactly to the
full-model $R^2=v(N)$, and we report each as a percentage share. As companion
descriptors we report each group's single-domain skill $v(\{i\})$ and the
\emph{overlap} $\sum_i v(\{i\}) - v(N)$, which quantifies how much the groups'
explanatory contributions coincide: a large overlap signifies predictors that are
correlated but not redundant, and --- following the argument of
Sec.~\ref{sec:intro} --- is itself an expression of the zonal coupling between
climate and soil rather than a defect of the design.

The biological group is carried in two forms. The MAIN configuration uses six
layers (soil fungal proportion and five SPUN mycorrhizal
richness/endemism/ratio layers; $\geq96\%$ coverage of modelling rows); the three
\citet{Barcelo2023} root-colonisation layers, which add negligible unique skill
($\Delta R^2\approx0.001$) while reducing coverage to 61\%, are retained only as a
nine-layer sensitivity case (Sec.~\ref{sec:robust}, Appendix).

% ---------------------------------------------------------------------------
\subsection{Spatial organisation of the beyond-climate residual (Moran's $I$)}
\label{sec:moran}
% ---------------------------------------------------------------------------
\textit{In brief.} Attribution says how much non-climate variance exists; it does
not say whether that variance is spatially \emph{organised} or merely scattered.
We test this with Moran's $I$, which asks whether places near one another carry
more similar values than they would if the values were shuffled at random across
space. $I\approx0$ indicates no spatial pattern; $I>0$ indicates that nearby
locations resemble each other --- the signature of an organised geography.

\textit{Formally.} We take the beyond-climate residual as the out-of-sample
$\ln\tau$ minus the climate-only prediction, aggregate it to block means, and
compute
$I=\frac{n}{\sum_{ij}w_{ij}}\cdot
\frac{\sum_{ij}w_{ij}(x_i-\bar x)(x_j-\bar x)}{\sum_i (x_i-\bar x)^2}$,
where $w_{ij}$ is a spatial contiguity weight between blocks $i$ and $j$.
Significance is assessed against a null distribution from 999 random spatial
permutations of the residuals. We evaluate $I$ at $2\textdegree$ and
$5\textdegree$.

% ---------------------------------------------------------------------------
\subsection{Mapping the zonality, and choosing the scale}
\label{sec:mapscale}
% ---------------------------------------------------------------------------
\textit{In brief.} To show \emph{where} non-climate context bends turnover, we map
the ratio of a context-aware prediction to a climate-only prediction: values above
one mark regions where edaphic or biological context lengthens turnover relative to
what climate alone implies, values below one where it shortens it. Such a map is
only worth drawing at a scale where it is reliable. Two analyses set that scale.
First, most of the point-to-point variation in the residual is unstructured: a
variogram shows roughly two-thirds of its variance is short-range noise. Second,
and more directly, we measured how faithfully the map reproduces the observed
residual as it is aggregated to coarser blocks; agreement roughly doubles from the
noisy native pixel to a meso scale and then plateaus, with no single privileged
resolution. We therefore map at $2\textdegree$, within that plateau.

\textit{Formally.} The modulation field is the log-ratio of model predictions,
$\ln[\hat\tau_{\text{context}}/\hat\tau_{\text{climate}}]$, computed symmetrically
(each domain relative to climate, with no abiotic-before-biology ordering) to match
the order-independent attribution; the abiotic and biological panels are shown
separately. The noise floor is characterised by an empirical variogram of the
residual (nugget/sill ratio and range). The mapping scale is set by an
out-of-sample skill curve: using spatial-block-CV predictions, we form the native
beyond-climate map $\hat r_i=\hat\tau_{\text{full}}-\hat\tau_{\text{climate}}$ and
the observed residual $r_i$, aggregate both to block means at a series of block
sizes, and report $R^2(\bar r,\hat{\bar r})$ as a function of scale. This agreement
rises from $\approx0.11$ at the native pixel to $\approx0.22$ at
1--2\textdegree{} and is flat across 1--7\textdegree; a purely linear version of the
same diagnostic instead declines with scale, so the meso-scale choice rests on the
nonlinear, map-faithful metric (Appendix).

% ---------------------------------------------------------------------------
\subsection{Mechanism: marginal responses (ALE)}
\label{sec:ale}
% ---------------------------------------------------------------------------
\textit{In brief.} Attribution and the map say how much and where; to say \emph{how}
each driver acts we trace the response of $\tau$ to each predictor while letting the
other predictors keep their naturally co-occurring values. Accumulated Local
Effects (ALE) do this without the extrapolation artefacts that partial-dependence
curves suffer when predictors are correlated.

\textit{Formally.} We use ALE \citep{Apley2020} on the MAIN 6-layer model
(consistent with the attribution and the maps). Each continuous predictor is split
into 40 quantile bins; within each bin the focal variable is nudged to the bin
edges with all others held at observed values, and the local prediction differences
are accumulated and centred. Curves are computed globally and within four latitude
bands (southern extratropics, tropics, northern temperate, boreal/Arctic).
Categorical predictors (K\"oppen zone, soil functional group) are excluded.

% ---------------------------------------------------------------------------
\subsection{Comparison with Earth System Model structure}
\label{sec:esm}
% ---------------------------------------------------------------------------
We compare the spatial structure of our $\tau$ against six CMIP6 models spanning
structurally distinct land-surface lineages \citep{Varney2022}: CESM2, UKESM1-0-LL,
IPSL-CM6A-LR, MPI-ESM1-2-LR, CanESM5, MIROC-ES2L. For each we computed the
quasi-equilibrium turnover $\tau_\text{ESM}=\text{cSoil}/\text{rh}$
\citep{Koven2017,Carvalhais2014} over 1980--2014 (using \texttt{cSoilAbove1m} for
CESM2 to reduce depth mismatch).

Absolute turnover is not comparable across products: depth, pool definition, and
reference period all differ, and these differences are unavoidable in cross-ESM
$\tau$ comparison. We therefore do not treat any product as a standard of truth and
do not certify absolute skill; we compare only the normalised spatial
\emph{structure}. Each product is normalised to its own distribution, and structure
is summarised two ways: relative turnover by biome zone (main text) and by
continuous latitude band (appendix), together with magnitude-free dispersion
statistics (coefficient of variation, 95:5 percentile ratio, SD of $\ln\tau$). We
quantify over-amplification as the ratio of polar to tropical relative turnover in
each product against the same ratio in the observation-constrained map.

% ---------------------------------------------------------------------------
\subsection{Robustness checks}
\label{sec:robust}
% ---------------------------------------------------------------------------
Three checks bound the main claims (details and figures in the Appendix).
\emph{Provenance.} To test whether the residual structure could be an artefact of
combining many source databases, we measured the between-source share of residual
variance ($\eta^2$) and the change in coarse organisation after removing each
source's mean. \emph{Biological parsimony.} We compared per-layer coverage,
collinearity, and unique skill to justify the MAIN 6-layer set against the
9-layer alternative. \emph{Block-size sensitivity.} We repeated the Shapley
decomposition at $1$, $2$, and $5\textdegree$ blocks (Sec.~\ref{sec:rf}).

% ---------------------------------------------------------------------------
\subsection{Software and reproducibility}
% ---------------------------------------------------------------------------
Analyses used R~($\geq4.3$) with \texttt{terra}, \texttt{sf}, \texttt{ranger}
\citep{Wright2017}, and \texttt{gstat} (variograms). The pipeline is a sequence of
numbered scripts from data assembly through prediction, attribution, organisation,
mapping-scale, sensitivity, and ESM comparison; each reported number and figure
regenerates from a named script. Code and intermediate outputs are archived at
[DOI].
